\documentclass[12pt]{article}
\usepackage[utf8]{inputenc}
\usepackage[T1]{fontenc}
\usepackage{hyperref}
\title{Tuscarora Language Resources for K-12 Curriculum Development}
\author{Compiled by Claude Sonnet 5 (claude-sonnet-5) for Andr\'as Kornai}
\date{}
\begin{document}
\maketitle

\section{Introduction}
Tuscarora (ISO 639-3 tus) is an Iroquoian language of the Tuscarora Nation in New York and of Six Nations in Canada. It is dormant, with very few speakers. The most useful published reference is Blair Rudes's dictionary. Search results also point to a revitalization programme and a workbook, but I could not confirm current programme details.

\section{Primary Educational Resources}
\begin{itemize}
\item \emph{Daughters, Sisters, Mothers \& Wives: A Tuscarora Language Workbook}, with short lessons including the Thanksgiving Address and self-introduction. \url{https://onlc.ca/bookstore/p/daughters-sisters-mothers-wives-a-tuscarora-language-workbook}
\end{itemize}

\section{Linguistic and Reference Materials}
\begin{itemize}
\item Rudes, \emph{Tuscarora-English / English-Tuscarora Dictionary} (University of Toronto Press), paperback listing. \url{https://www.goodminds.com/collections/all-products/products/tuscarora-english-english-tuscarora-dictionary-pb}
\item Library record for the same dictionary. \url{https://openlibrary.org/books/OL117971M}
\item Native Languages of the Americas: Tuscarora vocabulary and pronunciation pages. \url{https://www.native-languages.org/tuscarora.htm}
\item Endangered Languages Project entry for Tuscarora. \url{https://www.endangeredlanguages.com/node/109764}
\end{itemize}

\section{Religious and Community Resources}
The Tuscarora Language Revitalization Program (TLRP) in Lewiston, NY is described as running adult classes, immersion and master-apprentice sessions. See the talking dictionary project page below. I did not find the programme's own website.

\section{Digital and Audio Resources}
\begin{itemize}
\item Community-driven Tuscarora talking dictionary (2013 Global Voices microgrant project). I could not confirm whether it was completed. \url{https://rising.globalvoices.org/microgrants2013/usa-community-driven-tuscarora-talking-dictionary/}
\end{itemize}

\section{Curriculum Development Context}
Tuscarora is not known to be taught in schools. The dictionary, the workbook and any TLRP audio are the likeliest starting points, and the Tuscarora community should decide how they are used. Dictionary prices vary widely (one source warns of very high resale prices), so check with libraries. Contact the TLRP and the Tuscarora Nation directly about current materials.
\end{document}